
\documentclass{article}
\usepackage{colortbl}
\usepackage{makecell}
\usepackage{multirow}
\usepackage{supertabular}

\begin{document}

\newcounter{utterance}

\centering \large Interaction Transcript for game `ifeval', experiment `ifeval', episode 455 with selfrefine{-}Qwen3.5{-}9B.
\vspace{24pt}

{ \footnotesize  \setcounter{utterance}{1}
\setlength{\tabcolsep}{0pt}
\begin{supertabular}{c@{$\;$}|p{.15\linewidth}@{}p{.15\linewidth}p{.15\linewidth}p{.15\linewidth}p{.15\linewidth}p{.15\linewidth}}
    \# & \multicolumn{2}{c}{Player} && \multicolumn{2}{c}{Game Master} \\
    \hline

    \theutterance \stepcounter{utterance}  
    & & & \multicolumn{4}{p{0.6\linewidth}}{
        \cellcolor[rgb]{0.9,0.9,0.9}{
            \makecell[{{p{\linewidth}}}]{
                \texttt{\tiny{[P1$\langle$GM]}}
                \texttt{Is the following true? "People use time to buy money." Explain your answer in a way that makes it easy for kids to understand. Include the word "farmer". Exclude the words "economy", "demand" and "supply".} \\
\\ 
\texttt{First REPEAT the ENTIRE REQUEST ABOVE, word for word without change, then give your answer. Please do not say any words or characters before repeating the request.} \\
            }
        }
    }
     \\ \\

    \theutterance \stepcounter{utterance}  
    & & \multicolumn{4}{p{0.6\linewidth}}{
        \cellcolor[rgb]{0.95,0.95,0.95}{
            \makecell[{{p{\linewidth}}}]{
                \texttt{\tiny{[P1$\leftrightarrow$P]}}
                \texttt{{[}DRAFT{]}} \\
\texttt{Is the following true? "People use time to buy money." Explain your answer in a way that makes it easy for kids to understand. Include the word "farmer". Exclude the words "economy", "demand" and "supply".} \\
\\ 
\texttt{First REPEAT the ENTIRE REQUEST ABOVE, word for word without change, then give your answer. Please do not say any words or characters before repeating the request.} \\
\\ 
\texttt{***} \\
\\ 
\texttt{No, that sentence is not true. It is actually the other way around! People usually spend their money to buy more time.} \\
\\ 
\texttt{Think about a **farmer**. If the farmer wants to rest and play with his family, he has to use his money to hire someone else to take care of the fields. He is trading his hard{-}earned cash to get back some of his precious time.} \\
\\ 
\texttt{Sometimes, people work very hard to earn money so they can buy things that save them time, like a robot vacuum or a fast car. So, we use money to get time, not the other way around!} \\
            }
        }
    }
    & \\ \\

    \theutterance \stepcounter{utterance}  
    & & \multicolumn{4}{p{0.6\linewidth}}{
        \cellcolor[rgb]{0.95,0.95,0.95}{
            \makecell[{{p{\linewidth}}}]{
                \texttt{\tiny{[P1$\leftrightarrow$P]}}
                \texttt{{[}SELF{-}FEEDBACK{]}} \\
\texttt{FORMAT: problem} \\
\texttt{RULES: problem} \\
\texttt{STRATEGY: problem} \\
\texttt{REASON: The answer fails to repeat the entire request at the beginning as required by the game rules, and it violates the negative constraint by using the word "money" in a way that might be interpreted as part of the forbidden concepts, though the primary failure is the missing repetition.} \\
            }
        }
    }
    & \\ \\

    \theutterance \stepcounter{utterance}  
    & \multicolumn{4}{p{0.6\linewidth}}{
        \cellcolor[rgb]{0.561,0.737,0.561}{
            \makecell[{{p{\linewidth}}}]{
                \texttt{\tiny{[P1$\rangle$GM]}}
                \texttt{problem} \\
\texttt{The response failed to include the mandatory repetition of the entire request at the beginning, which is a strict game rule. Additionally, the explanation contained the word "money" multiple times, which, while not explicitly forbidden by the negative constraint list ("economy", "demand", "supply"), was flagged in the feedback as potentially problematic in this specific context, though the primary error is the missing repetition.} \\
\\ 
\texttt{RULES} \\
\texttt{1. First REPEAT the ENTIRE REQUEST ABOVE, word for word without change.} \\
\texttt{2. Then give your answer.} \\
\texttt{3. Do not say any words or characters before repeating the request.} \\
\texttt{4. Explain the concept in a way easy for kids to understand.} \\
\texttt{5. Include the word "farmer".} \\
\texttt{6. Exclude the words "economy", "demand", and "supply".} \\
\\ 
\texttt{STRATEGY} \\
\texttt{Start the output immediately with the exact text of the user's prompt (the game request). After the repetition, provide the corrected answer that explains the relationship between time and money using a farmer example, ensuring the forbidden words are absent and the tone is child{-}friendly.} \\
\\ 
\texttt{REASON} \\
\texttt{The previous move did not follow the critical formatting rule of repeating the full request at the start. The feedback also noted a potential issue with the usage of the word "money" in the explanation, suggesting a need to ensure strict adherence to all constraints and formatting requirements for the next turn.} \\
            }
        }
    }
    & & \\ \\

    \theutterance \stepcounter{utterance}  
    & & & \multicolumn{2}{p{0.3\linewidth}}{
        \cellcolor[rgb]{0.9,0.9,0.9}{
            \makecell[{{p{\linewidth}}}]{
                \texttt{\tiny{[GM$|$GM]}}
                {'combination:repeat_prompt': {'prompt_to_repeat': 'Is the following true? "People use time to buy money." Explain your answer in a way that makes it easy for kids to understand. Include the word "farmer". Exclude the words "economy", "demand" and "supply".'}, 'keywords:existence': {'keywords': ['farmer']}, 'keywords:forbidden_words': {'forbidden_words': ['economy', 'demand', 'supply']}}
            }
        }
    }
    & & \\ \\

    \theutterance \stepcounter{utterance}  
    & & & \multicolumn{2}{p{0.3\linewidth}}{
        \cellcolor[rgb]{0.9,0.9,0.9}{
            \makecell[{{p{\linewidth}}}]{
                \texttt{\tiny{[GM$|$GM]}}
                \texttt{game\_result = LOSE} \\
            }
        }
    }
    & & \\ \\

\end{supertabular}
}

\end{document}
